\chapter{Conclusion}
\label{chap:conclusion}

In 2018, the four of us built, in a few weeks, a segmentation that the researchers approved, and a
viewer for it. We had identified the real difficulty, telling vessels from bone, and answered it with
morphology tuned by eye.

Eight years later, pyBones keeps the goal and changes the method, and replaces our viewer with a local
web page that segments a stack, shows its slices, turns its meshes around in 3D, measures its
vessel network and compares two runs: the quantification the brief asked for, with its limits stated. The main lessons I
draw from the rewrite are:
\begin{itemize}
  \item \textbf{Measure before improving.} Without a ground truth, a synthetic phantom calibrated on the
        real scan made every change measurable, and showed me that an intuitive improvement (the first 3D
        refinement) made things worse until the vessel edges were right.
  \item \textbf{Physics over parameters.} The largest gains came from modelling how the image forms
        (partial volume, blur, density) rather than from tuning radii: half-maximum edges for vessels
        and bursts, physical units for every size.
  \item \textbf{3D where 2D cannot decide.} Connectivity and shape in 3D resolve most of the blobs that a
        single slice leaves ambiguous; what remains is a property of the scan (resolution) and of the
        anatomy (rods are tubes too). A density measured in one slice even mistook a vessel lying in it
        for the edge of a burst.
\end{itemize}

\section*{Future work}
\addcontentsline{toc}{section}{Future work}

\begin{itemize}
  \item \textbf{Expert labels}: a few real crops annotated with the annotation page would validate
        the segmentation on the real scan, and decide the two experimental switches (the bone-side
        vessel edges and the thin vessels), which improve the phantoms but stay off until then
        (Section~\ref{sec:experimental}).
  \item \textbf{Cleaner junctions}: keeping clotted segments joined to their vessels, and telling
        apart vessels the blur fuses. The defects of the vessel mask,
        not the measure, make the errors of the network (Chapter~\ref{chap:network}).
  \item \textbf{Measures per region} (cortex, marrow), to compare the network across the bone.
  \item \textbf{Gap bridging} along the vessel direction, to reconnect vessels interrupted by clots.
  \item \textbf{The web page}: the raw slice without overlay, and measurements in 3D (distances,
        the calibre of a picked vessel).
\end{itemize}
